% Round 15. New candidate following the independently reviewed incidence
% formula. Requires a separate sign and boundary-category audit.
\section{The incidence homotopy on the outer edges and the first strip}
\label{r15s:section}

This section keeps the incidence geometry, smooth line curvatures,
scalar $g$, and polynomial $W$ from
Section~\ref{r15i:section}.  Thus $E=\mathbf P_S(N)$ has
relative dimension $r\geq1$, $F=\mathbf P_E(Q)$,
$I=\mathbf P_F(M)$, and $p=\lambda\rho:I\to E$.
Write $m=\dim E$ and $t=\pi\lambda:F\to S$.
In addition to $(\mathrm{HP}_{\mathrm{nat}})_K$, the proofs
below use the separately proved PWS target projection formula
\eqref{r15ps:PWS-projection}.  This is an explicit
tensor-level result, with its own band-integration proof;
it is not treated as a formal extension of smooth testing.

All Aeppli groups below are actual groups
\[
 H_{\mathrm A}^{a,b}(V)
 =\ker(D:B^{a,b}(V)\to B^{a+1,b+1}(V))
     /\bigl(d'B^{a-1,b}(V)+d''B^{a,b-1}(V)\bigr).
\]
Negative bidegrees are zero.  In particular
$H_{\mathrm A}^{a,0}=\ker D/d'B^{a-1,0}$.
Every actual form on an outer edge is PWS.

\begin{lemma}[The PWS incidence homotopy in every bidegree]
\label{r15s:general-homotopy}
Let $\alpha\in\operatorname{PS}^{a,b}(E)$ satisfy $D\alpha=0$.
Set
\[
 \beta_\alpha=Wz^*\alpha,
 \quad C\alpha=\rho_*(z^*\alpha),
 \quad J\alpha=\pi_*(\xi^r\alpha).
\tag{FS15.1}\label{r15s:CJ}
\]
There are ordinary PWS outputs
\[
 \mathcal G\alpha\in\operatorname{PS}^{a-1,b-1}(E),\quad
 \mathcal P\alpha\in\operatorname{PS}^{a,b-1}(E),\quad
 \mathcal Q\alpha\in\operatorname{PS}^{a-1,b}(E)
\tag{FS15.2}\label{r15s:degrees}
\]
representing the prescribed singular expressions with
$g\beta_\alpha,d'g\beta_\alpha,d''g\beta_\alpha$,
respectively.  The outputs of negative degree are zero.
The following is an equality of ordinary actual forms:
\[
 \alpha=\pi^*(J\alpha)
        -\lambda_*(\nu^r\wedge C\alpha)
        +D\mathcal G\alpha+d''\mathcal P\alpha
                                      -d'\mathcal Q\alpha.
\tag{FS15.3}\label{r15s:homotopy}
\]
Both $C\alpha$ and $J\alpha$ are $D$-closed PWS forms,
of degrees $(a-1,b-1)$ and $(a,b)$, respectively.
\end{lemma}

\begin{proof}
The germ identity $z^*\alpha=p^*\alpha$ on an ordinary
neighborhood of $\Sigma$ holds in every bidegree.  For the
tests complementary to the three degrees in
\eqref{r15s:degrees}, multiplication by $\alpha$ gives
respectively target degrees $(m+1,m+1),(m,m+1),(m+1,m)$.
Thus the same cutoff argument as in
Proposition~\ref{r15i:outputs} produces global PWS inputs
for the projective-bundle tower and proves every prescribed
smooth-test pairing.  The input degrees are
$(a+r-1,b+r-1),(a+r,b+r-1),(a+r-1,b+r)$;
subtracting the relative degree $(r,r)$ gives exactly
\eqref{r15s:degrees}.  There is no assertion about an
operator on forms of positive delta codimension.

For clarity we redo the signs, rather than extrapolate the
even-degree calculation.  Let
$\eta\in\mathcal A_c^{m-a,m-b}(E)$ and write
$\eta_p=p^*\eta$ and $\epsilon=(-1)^{a+b}$.
The total degree of $\beta_\alpha$ has parity $a+b$.
Put
\[
 A=\int_I g\beta_\alpha D\eta_p,\quad
 B=\int_I d'g\,\beta_\alpha d''\eta_p,\quad
 C=\int_I d''g\,\beta_\alpha d'\eta_p.
\]
All displayed integrands vanish near $\Sigma$ whenever
needed; hence these are ordinary integrals of forms on the
complement, with compact support there.  The differential
convention for currents gives
\[
 \langle D\mathcal G\alpha+d''\mathcal P\alpha
                   -d'\mathcal Q\alpha,\eta\rangle
                       =A+\epsilon B-\epsilon C.
\tag{FS15.4}\label{r15s:current-sign}
\]
Actual Stokes for the zero extensions of
$g\beta_\alpha d''\eta_p$ and
$g\beta_\alpha d'\eta_p$ gives
\[
 \int_I g\,d'\beta_\alpha d''\eta_p=-B-\epsilon A,
 \qquad
 \int_I g\,d''\beta_\alpha d'\eta_p=-C+\epsilon A.
\]
Since $D\beta_\alpha=0$,
\[
 D(\beta_\alpha\eta_p)
   =-\epsilon d''\beta_\alpha d'\eta_p
      +\epsilon d'\beta_\alpha d''\eta_p
      +\beta_\alpha D\eta_p.
\]
It follows that the right side of
\eqref{r15s:current-sign} is
$-\int_IgD(\beta_\alpha\eta_p)$.
This also handles the outer edges: the inadmissible tests
and corresponding integrals are zero by degree.

Apply the scalar Green formula to the actual test
$\beta_\alpha\eta_p$, of degree
$(m+r-1,m+r-1)$.  On $\Sigma$, terms in $W$ with
$k<r-1$ vanish because $\alpha\eta$ is a target top form.
The remaining term integrates to $\int_E\alpha\eta$
by the frozen projective fiber-degree formula (R10.1),
valid against every actual top test.  The term containing
$(z^*\xi)^r$ is $\int_E\pi^*(J\alpha)\eta$ by the
top-form change of variables away from the diagonal and
projective-bundle base change, exactly as in
\eqref{r15i:first-term}.  The top form restricts to zero
on the diagonal by dimension, independently of $(a,b)$.
For the term containing $\nu^r$, first test $\rho_*z^*\alpha$
against the smooth form $\nu^r\lambda^*\eta$, and then
test $\lambda_*(\nu^r C\alpha)$ against $\eta$.
This gives \eqref{r15s:homotopy} for all smooth tests,
and actual-current injectivity proves it as stated.
Closedness of $C\alpha,J\alpha$ follows by current
differentiation and actual-current injectivity.
\end{proof}

\begin{theorem}[Projective bundles do not add outer-edge Aeppli classes]
\label{r15s:edge}
For every $a\geq0$,
\[
 H_{\mathrm A}^{a,0}(E)=\pi^*H_{\mathrm A}^{a,0}(S),
 \qquad
 H_{\mathrm A}^{0,a}(E)=\pi^*H_{\mathrm A}^{0,a}(S),
\tag{FS15.5}\label{r15s:edge-formula}
\]
with injective pullbacks.  More precisely, for a $D$-closed
actual form $\alpha$ of degree $(a,0)$ there are $D$-closed
$\gamma\in\operatorname{PS}^{a,0}(S)$ and
$q_\alpha\in\operatorname{PS}^{a-1,0}(E)$ such that
\[
                    \alpha=\pi^*\gamma-d'q_\alpha.
\tag{FS15.6}\label{r15s:edge-identity}
\]
The conjugate statement holds on the other edge.
\end{theorem}

\begin{proof}
The input is already PWS.  With $b=0$,
$C\alpha$, $\mathcal G\alpha$, and $\mathcal P\alpha$
have negative second degree, hence vanish.
Equation~\eqref{r15s:homotopy} becomes
\eqref{r15s:edge-identity}, with $\gamma=J\alpha$ and
$q_\alpha=\mathcal Q\alpha$.
Injectivity is the frozen constant-fiber Aeppli theorem,
Round 14 \texttt{r14j:equidimensional}, in minimum degree
zero.  The involution exchanging the two actual differentials
gives the second assertion.  A projective bundle of relative
dimension zero is an isomorphism, for which the same claim
holds with $q_\alpha=0$.
\end{proof}

\begin{lemma}[The line contraction respects first-strip Aeppli boundaries]
\label{r15s:C-boundaries}
Let $a\geq1$ and let
$\alpha\in\operatorname{PS}^{a,1}(E)$ be an actual
Aeppli boundary.  Then
\[
           C\alpha=d'\delta
           \quad\hbox{for some }\quad
           \delta\in\operatorname{PS}^{a-2,0}(F).
\tag{FS15.7}\label{r15s:C-boundary}
\]
For $a=1$ this says $C\alpha=0$.
\end{lemma}

\begin{proof}
Write $\alpha=d'\mu+d''\tau$, with
$\mu\in B^{a-1,1}(E)$ and $\tau\in B^{a,0}(E)$.
The second form is PWS.  If $a\geq2$, the frozen actual
first-strip replacement gives
$\mu=\mu_0+d'\mu_1$ with both $\mu_0,\mu_1$ PWS.
If $a=1$, $\mu$ is already on an outer edge and put
$\mu_0=\mu$.  In either case
$\alpha=d'\mu_0+d''\tau$.

The proper current pushforward of $z^*\mu_0$ under
$\rho$ has an ordinary PWS representative of degree
$(a-2,0)$, whereas that of $z^*\tau$ has negative second
degree and is zero.  Proper differentiation of currents
therefore proves
\[
        [C\alpha]=d'[\rho_*z^*\mu_0].
\]
Both sides are currents of actual forms; injectivity proves
\eqref{r15s:C-boundary}.  No form-valued pushforward of
$d'\mu_0$ or $d''\tau$ is used.  These derivatives may
indeed have positive delta codimension, and their current
identity is used only to differentiate already constructed
outputs.
\end{proof}

\begin{theorem}[The complete actual first-strip formula]
\label{r15s:first-strip}
Under the analytic inputs stated at the start of this section
and the frozen first-strip PWS replacement, for every
$a\geq1$ one has
\[
 H_{\mathrm A}^{a,1}(E)
  =\pi^*H_{\mathrm A}^{a,1}(S)
      \oplus [\xi]\wedge\pi^*H_{\mathrm A}^{a-1,0}(S).
\tag{FS15.8}\label{r15s:first-strip-formula}
\]
There is the conjugate formula
\[
 H_{\mathrm A}^{1,a}(E)
  =\pi^*H_{\mathrm A}^{1,a}(S)
      \oplus [\xi]\wedge\pi^*H_{\mathrm A}^{0,a-1}(S).
\tag{FS15.9}\label{r15s:conjugate-strip}
\]
The products are well defined because $\xi$ is separately
closed.  The formula uses the actual Bott--Chern class of
$\mathcal O_E(1)$; it does not depend on the chosen smooth
metric.
\end{theorem}

\begin{proof}
An arbitrary actual class of degree $(a,1)$ has a
$D$-closed PWS representative $\alpha$ by Round 14's
first-strip replacement.  The PWS form
$C\alpha\in\operatorname{PS}^{a-1,0}(F)$ is $D$-closed.
Apply the edge formula first to $\lambda:F\to E$ and
then to $\pi:E\to S$.  It gives a $D$-closed PWS form
$\gamma$ of degree $(a-1,0)$ on $S$ and a PWS form
$\delta$ of degree $(a-2,0)$ on $F$ with
\[
                       C\alpha=t^*\gamma+d'\delta.
\tag{FS15.10}\label{r15s:C-decomposition}
\]
For $a=1$ the last term is zero.  The relative-dimension-zero
case of $\lambda$ when $r=1$ causes no change.

The PWS target projection formula and the relative top
Chern normalization give
\[
       C(\xi\wedge\pi^*\gamma)
          =\rho_*(z^*\xi\wedge\rho^*t^*\gamma)
          =t^*\gamma.
\tag{FS15.11}\label{r15s:C-hyperplane}
\]
Put $\alpha_0=\alpha-\xi\pi^*\gamma$.  It is PWS and
$D$-closed, and $C\alpha_0=d'\delta$.
Since $\delta$ is on an outer edge, both $\delta$ and
$d'\delta$ are PWS.  The form
$\Lambda=\lambda_*(\nu^r\wedge\delta)$ is a PWS
form of degree $(a-1,1)$ on $E$, and current Stokes gives
the actual equality
\[
                  \lambda_*(\nu^r d'\delta)=d'\Lambda.
\tag{FS15.12}\label{r15s:boundary-push}
\]
Here both inputs to the PWS operator are legitimately PWS;
equivalently the identity follows by differentiating its
prescribed currents and then using current injectivity.

Apply \eqref{r15s:homotopy} to $\alpha_0$ and use
\eqref{r15s:boundary-push}.  It expresses $\alpha_0$ as
$\pi^*(J\alpha_0)$ plus actual Aeppli boundaries.
The base form $J\alpha_0$ is $D$-closed of degree $(a,1)$.
Together with $\alpha=\alpha_0+\xi\pi^*\gamma$ this
proves spanning in \eqref{r15s:first-strip-formula}.

For directness, represent the base classes by $D$-closed
PWS forms $\beta$ of degree $(a,1)$ and $\gamma$ of
degree $(a-1,0)$.  Suppose
$\pi^*\beta+\xi\pi^*\gamma$ is an actual Aeppli boundary.
Apply Lemma~\ref{r15s:C-boundaries}.  The line contraction
of $\pi^*\beta$ is zero by relative degree and the PWS
target projection formula.  The contraction of the second
term is $t^*\gamma$ by \eqref{r15s:C-hyperplane}.
Thus $[t^*\gamma]=0$ in $H_{\mathrm A}^{a-1,0}(F)$.
The injective edge pullbacks for $\pi$ and $\lambda$
give $[\gamma]=0$.  The second summand was therefore
already an Aeppli boundary, so $\pi^*[\beta]=0$.
The frozen constant-fiber Aeppli injection in minimum
degree one gives $[\beta]=0$.

If the metric is changed by a global PWS scalar $u$, the
curvature changes by $Du$.  For any $D$-closed actual
$\gamma$ the explicit identity
\[
       (Du)\wedge\gamma
          =d'(d''u\wedge\gamma)+d''(u\,d'\gamma)
\tag{FS15.13}\label{r15s:metric-independence}
\]
follows by the Leibniz rule and $D\gamma=0$.
Hence its Aeppli action is unchanged.  Interchanging the
two differentials proves \eqref{r15s:conjugate-strip}.
\end{proof}

\paragraph{Exact scope.}
The general homotopy lemma is a statement about $D$-closed
PWS inputs in every bidegree.  The complete actual cohomology
conclusions above use PWS replacement only on the outer
edges and the first strip, where the frozen theorem supplies
it.  No replacement of arbitrary actual $(2,2)$ forms, no
all-actual-test nonfinite pushforward, and no computation of
the higher actual Aeppli groups is inferred.
